\documentclass[a4paper, titlepage]{article}
\usepackage{kotex}
\usepackage{tabularx}
\usepackage[a4paper, 
            top=3cm,
            bottom=3cm,
            left=4.5cm,
            right=4.5cm]{geometry}

\begin{document}

\title{essential\_sample.\TeX}
\author{SeungHeon LEE}
\date{\today}
\maketitle

\tableofcontents

\newpage
\section{본문 \emph{Running Text}}
\LaTeX \ 은 입력파일에서 공백주거나 단어를 띄워두는 것이 출력물에 영향을 주지 않는다. 입력파일에서 공백이 몇 개이든지간에, \LaTeX \ 은 이를 하나의 공백으로 처리한다. 
또한, \LaTeX \ 은 각 행의 끝을 단어 간의 공백으로 간주하여 자동으로 띄어쓰기를 한다. 
입력파일에서 공백인 행이 존재할 경우 이는 새로운 단락의 시작임을 가리키므로, 새로운 단락을 시작하는 것이 아니라면 공백인 행을 두지 말자. 

\LaTeX \ 은 잘 쓰이지 않는 기호 
\begin{quote}
\#, \$, \%, \&, \_, \{, \}, \textasciitilde, \textasciicircum
\end{quote}
를 입력파일에 그냥 입력하면 제대로 처리하지 못한다.

\subsection{\LaTeX \ 명령어}
\begin{itemize}
    \item \LaTeX 에서 사용되는 명령어는 항상 `\textbackslash' 로 시작한다.  
    \item \LaTeX 명령어는 대소문자를 구분한다는 점을 유의하여야 한다.
    \item 다음과 같은 형태의 명령어도 있다. 
    \begin{quote}
        \textbackslash command\{text\}  
    \end{quote}
    이때, \{ \} 안에 들어가는 text는 명령어의 인자(parameter 또는 argument)이며, 필수적으로 들어가야 한다. 
    반면, [ ] 안에 들어가는 text는 없어도 되는 선택 인자이다. 
\end{itemize}

\subsection{기타 살펴볼 사항}
\begin{itemize}
    \item \LaTeX \ 은 열고 닫는 홑따옴표와 겹따옴표를 입력할 수 있다. 여는 따옴표는 `를 활용하고, 닫는 따옴표는 '를 활용한다. 겹따옴표는 앞의 따옴표를 두 번 입력하면 된다. 
    \item \LaTeX \ 에는 총 세 가지의 대시가 존재한다. 하이픈(-), 짧은 대시(--)\footnote{숫자의 범위 표시에 주로 사용. ex. \ 10--20}, 긴 대시(---)\footnote{문장의 구분에 사용.}가 그것이다.  
    \item 본문에 특수문자를 입력해야 할 때가 있다. 이때는 백슬래시(\textbackslash) \ 를 활용하여 입력할 수 있다. 
\end{itemize}

\section{도큐먼트 클래스 \emph{documentclass}}
본 절에서는 \LaTeX \ 에 존재하는 몇 가지 표준 도큐먼트 클래스에 대해 소개한다. 
\begin{itemize}
    \item [\textbf{article}] 짧은 문서나 출판물 속에 들어갈 기사와 같은 것에 어울린다. article은 chapter가 없고 \textbackslash maketitle로 제목을 작성하면, 표지를 별도로 사용하지 않고 첫 페이지의 상단에 제목을 인쇄한다.
    \item [\textbf{report}] 좀더 기술적인 문서 작성에 알맞다. 이는 chapter를 가지며 제목을 별도의 페이지에 인쇄하는 점을 제외하면 article과 같다.
    \item [\textbf{book}] 책자의 출판을 목적으로 한다. 종이의 양면에 인쇄할 것으로 가정하여 페이지의 구성을 조정한다.
    \item [\textbf{letter}] 편지 등을 작성하기 위한 것이다. 편지 형식을 사용하면 주소, 날짜, 서명 등의 각 구성요소가 적절하게 배치된 편지를 작성할 수 있다.
\end{itemize}

또한, 이러한 표준 형식을 바꿀 수 있는 여러 옵션이 존재한다. 문서에서 도큐먼트 클래스는 하나만 사용할 수 있지만, 옵션은 여러 개 사용할 수 있다는 점이 그 특징이다. \LaTeX \ 에 존재하는 표준 옵션들을 일부 알아보자.\footnote{다른 옵션에 관한 정보를 알고 싶다면 \LaTeX :Document Structure 페이지를 참고하라.}  
\begin{itemize}
    \item [\textbf{11pt, 12pt}] 본문에 사용되는 폰트 크기는 일반적으로 10pt이지만, 해당 옵션을 활용하여 폰트의 크기를 키울 수 있다.
    \item [\textbf{twoside}] article이나 report 형식에서 양면 인쇄를 하기 위한 옵션이며, book에서는 twoside가 기본 설정에 포함된다. 
    \item [\textbf{titlepage}] article 형식에서 \textbackslash maketitle 명령어가 새로운 페이지에 제목을 작성하도록 한다. 이는 report나 book 형식에서는 기본 설정에 포함된다. 
    \item [\textbf{a4paper}] 모든 표준 형식의 출력을 a4 용지 크기에 맞추도록 한다.  
\end{itemize}

\section{환경}
\subsection{인용문 \emph{quote, quotation}}
\begin{itemize}
    \item [\textbf{quote}] quote는 짧은 인용문이 쓰일 때 활용한다. 
US presidents have been
known for their pithy remarks:
\begin{quote}
The buck stops here.
I am not a crook.
\end{quote}

    \item [\textbf{quotation}] quotation은 한 단락을 넘어가는 긴 인용문이 쓰일 때 활용한다. quotation은 단락별로 들여쓰기가 가능하다.
Here is some advice to remember:
\begin{quotation}
Environments for making quotations
can be used for other things as well.
Many problems can be solved by
novel applications of existing
environments.
\end{quotation}

    \item [\textbf{verse}] verse는 시를 작성할 때 활용한다. 각 행은 \textbackslash \textbackslash 로 구분하며, 연은 한 줄을 띄어쓴다.
\begin{verse}
I will arise and go now,
and go to Innisfree,\\
And a small cabin build there,
of clay and wattles made:\\
Nine bean-rows will I have there,
a hive for the honey-bee;\\
And live alone in the bee-loud glade.

And I shall have some peace there,
for peace comes dropping slow, \\
Dropping from the veils of the morning
to where the cricket sings; \\
There midnight's all a glimmer,
and noon a purple glow, \\
And evening full of the linnet's wings.
\end{verse}
\end{itemize}

\subsection{중앙 및 좌우정렬 \emph{center, flushleft, flushright}}
\begin{center}
    center는 중앙정렬, \\
\end{center} 
\begin{flushleft}
    flushleft는 좌측정렬, \\
\end{flushleft}
\begin{flushright}
    flushright는 우측정렬을 의미한다.
\end{flushright}

\subsection{나열 문단 \emph{itemize, enumerate, description}}
리스트를 작성하기 위한 환경으로는 총 세 가지가 존재한다. 각각의 환경에서 새로운 항목은 \textbackslash item으로 시작한다. enumerate 환경은 각 항목을 넘버링하는 반면, itemize 환경에서는 특정한 마크로 표시된다. 이때, 특정한 마크는 설정할 수 있다.
설명형 리스트를 작성하기 위한 환경이 description이다. \textbackslash item 뒤의 각괄호에 항목의 레이블을 지정할 수 있다. 
\begin{itemize}
\item Itemized lists are handy. 
\item However, don't forget
    \begin{enumerate}
    \item The `item' command.
    \item The `end' command.
    \end{enumerate} 
\end{itemize}
\begin{description}
\item [What] This is description.
\item [How] type \textbackslash begin\{description\} 
\end{description}

\subsection{그대로 보이기 \emph{verbatim}}
경우에 따라서는 터미널 화면에 나타난 모양대로 본문을 입력할 필요가 있다. 
이러한 경우에는 verbatim 환경을 활용할 수 있다. 한두글자 또는 단어에 대하여 그대로 보이게 하려면 \textbackslash verb라는 명령어를 쓴다. 
\begin{verbatim}
{ this finds %a & %b }

for i := 1 to 27 do 
    begin 
    table[i] := fn(i);
    process(i)
    end;
\end{verbatim}
\textbackslash verb를 활용할 때는 특정한 문자를 시작과 끝문자로 사용해야 한다. 이때, *는 시작과 끝문자로 활용해서는 안된다. 
\\ \verb+&#$@+ → \textbackslash verb를 활용하면 이와 같은 문자 입력을 할 수 있다.

\section{글꼴}
\LaTeX{}에서 글꼴을 변경하기 위한 명령어의 모음은 아래와 같다.

\begin{center}
\begin{tabularx}{\textwidth}{llX}
\em Command & \em or & \em Effect \\
\textbackslash textrm\{\dots\} & \{\textbackslash rmfamily\dots\} & Text is set in \textrm{roman} family \\
\textbackslash textsf\{\dots\} & \{\textbackslash sffamily\dots\} & Text is set in \textsf{sans serif} family \\
\textbackslash texttt\{\dots\} & \{\textbackslash ttfamily\dots\} & Text is set in \texttt{typewriter} family \\
\textbackslash textmd\{\dots\} & \{\textbackslash mdseries\dots\} & Text is set in \textmd{medium} series \\
\textbackslash textbf\{\dots\} & \{\textbackslash bfseries\dots\} & Text is set in \textbf{bold} series \\
\textbackslash textup\{\dots\} & \{\textbackslash upshape\dots\} & Text is set in \textup{upright} shape \\
\textbackslash textit\{\dots\} & \{\textbackslash itshape\dots\} & Text is set in \textit{italic} shape \\
\textbackslash textsl\{\dots\} & \{\textbackslash slshape\dots\} & Text is set in \textsl{slanted} shape \\
\textbackslash textsc\{\dots\} & \{\textbackslash scshape\dots\} & Text is set in \textsc{small caps} shape \\
\textbackslash emph\{\dots\} & \{\textbackslash em\dots\} & Text is set \em emphasized \\
\textbackslash textnormal\{\dots\} & \{\textbackslash normalfont\dots\} & Text is set in \textnormal{the document} font \\
\end{tabularx}
\end{center}

한편, 글꼴의 크기를 변경하는 명령어도 존재한다. \\
\tiny tiny
\scriptsize scriptsize
\footnotesize footnotesize
\small small
\normalsize normalsize \\
\large large
\Large Large 
\LARGE LARGE
\huge huge
\Huge Huge 

\section{섹션 명령과 목차}
\normalsize
\LaTeX{}의 섹션 명령어에는 다음과 같은 것들이 있다.
\begin{center}
\textbackslash chapter, \textbackslash section, \textbackslash subsection, \textbackslash subsubsection, \textbackslash paragraph, \textbackslash subparagraph
\end{center}

대부분의 문서에서 \textbackslash paragraph와 \textbackslash subparagraph에서 만들어진 제목에는 번호를 부여하지 않는다. 
paragraph와 subparagraph는 단지 section의 하위 레벨일 뿐이다. 한편, article 형식에서는 \textbackslash chapter를 사용할 수 없다.
사용자는 문서에서 \textbackslash tableofcontents를 사용하여 각종 섹션 명령을 바탕으로 작성된 목차를 작성하고 포함할 수 있다. 

\section{특수기호의 식자}
사용자는 키보드에 없는 다양한 종류의 특수기호들을 문서에 나타낼 수 있다. 
먼저, \LaTeX{}에서는 어떠한 글자에도 악센트를 표시할 수 있다. 
\begin{center}
\`{o}, \'{o}, \~{o}, \"{o} 
\end{center}
기타 특수기호와 그 입력 명령어는 다음과 같다. 
\begin{center}
\copyright, \S, \P, \ldots, \LaTeX, \TeX
\end{center}
\textbackslash textbackslash, \textbackslash textasciitilde, \textbackslash textasciicircum 등과 같이 \textbackslash{}를 활용하여 특수기호를 식자할 경우에는, 명령어 뒤에 \{\}을 활용하여 공백을 입력할 수 있다.
`\LaTeX is wonderful'과 달리, `\LaTeX{} is wonderful'이라고 입력하는 것이 그 예다. 

\section{표의 사용}

\section{편지}
`letter\_sample.\TeX' 문서 참조

\section{에러}

\section{패키지}

\section{마치는 글}


\end{document}
